
\فصل{ویژگی‌های کلیدی نسخه‌ی اول}

\قسمت{ویژگی‌های اصلی}

این ویژگی‌ها هسته‌ی ساکنا را تشکیل می‌دهند و باید در نسخه‌ی اول ارائه شوند:

\زیرقسمت{سیستم احراز هویت و کنترل دسترسی}
ساکنا چهار نقش کاربری دارد: ساکن، مدیر ساختمان، کارکن خدماتی و مدیر سامانه.
هر نقش دسترسی‌های متفاوتی دارد و سیستم \lr{RBAC} این را تضمین می‌کند.
احراز هویت پیش‌نیاز هرگونه تعامل با سامانه است.

\زیرقسمت{مدیریت ساختمان و واحدها}
مدیر ساختمان می‌تواند اطلاعات کامل مجتمع را در سامانه ثبت کند:
تعداد طبقات، واحدها، ساکنین و کارکنان خدماتی.
هر ساکن اطلاعات واحد خود را در پروفایل می‌بیند.

\زیرقسمت{سیستم درخواست‌های خدماتی}
ساکنین درخواست تعمیرات یا خدمات ارسال می‌کنند.
مدیر آن‌ها را بررسی و به کارکن مناسب ارجاع می‌دهد.
کارکن پس از انجام کار، گزارش می‌دهد.
ساکن در هر لحظه می‌تواند وضعیت درخواست را ببیند.

\زیرقسمت{سیستم شارژ و کیف پول}
مدیر دوره‌های شارژ را تعریف می‌کند (ماهانه، فصلی).
ساکنین از طریق کیف پول داخلی شارژ پرداخت می‌کنند.
تاریخچه‌ی کامل پرداخت‌ها برای ساکن و مدیر قابل مشاهده است.
مدیر وضعیت بدهکاران را گزارش می‌گیرد.

\زیرقسمت{سیستم رزرو امکانات مشترک}
مدیر امکانات قابل رزرو (استخر، سالن، باشگاه، پارکینگ) را با ظرفیت و زمان‌بندی تعریف می‌کند.
ساکنین آنلاین رزرو می‌کنند و سیستم به‌صورت خودکار از تداخل جلوگیری می‌کند.
امکان لغو رزرو نیز وجود دارد.

\زیرقسمت{سیستم اطلاعیه}
مدیر ساختمان اطلاعیه‌های رسمی منتشر می‌کند.
ساکنین اطلاعیه‌ها را مشاهده می‌کنند و آرشیو کامل در دسترس است.

\زیرقسمت{داشبورد گزارش‌گیری}
مدیر ساختمان به داشبوردی با گزارش‌های زیر دسترسی دارد:
\شروع{فقرات}
\فقره وضعیت پرداخت شارژها و فهرست بدهکاران
\فقره درخواست‌های خدماتی (ثبت‌شده، در جریان، تکمیل‌شده)
\فقره رزروهای امکانات مشترک
\پایان{فقرات}

\قسمت{ویژگی‌های ثانویه}

این ویژگی‌ها در صورت وجود زمان کافی پیاده‌سازی می‌شوند:

\شروع{فقرات}
\فقره رزرو امکانات با هزینه یا رایگان، با تعیین بازه‌ی زمانی
\فقره کیف پول ساختمان و تسویه‌ی خودکار حق‌الزحمه‌ی کارکنان خدماتی
\فقره تقسیم هزینه‌ها با ضریب‌گذاری برای واحدها
\پایان{فقرات}

\قسمت{ویژگی‌های اختیاری}

این ویژگی‌ها در نسخه‌های بعدی در نظر گرفته شده‌اند:

\شروع{فقرات}
\فقره پیام‌رسانی داخلی بین ساکنین و مدیر
\فقره سیستم نظرسنجی ساختمانی
\فقره اتصال به درگاه پرداخت اینترنتی واقعی
\فقره اعلان‌های \lr{SMS} و ایمیل
\فقره اپلیکیشن موبایل
\پایان{فقرات}


\قسمت{نیازمندی‌های غیرعملکردی (\lr{Non-Functional Requirements})}

جدول~\ref{جدول:NFR} نیازمندی‌های غیرعملکردی ساکنا را خلاصه می‌کند.

\begin{table}[ht]
\centering
\caption{نیازمندی‌های غیرعملکردی}
\label{جدول:NFR}
\begin{tabular}{|>{\raggedleft\arraybackslash}p{3cm}|>{\raggedleft\arraybackslash}p{9cm}|}
\hline
\rl{\textbf{دسته}} & \rl{\textbf{نیازمندی}} \\
\hline
\rl{امنیت} & \rl{احراز هویت الزامی برای دسترسی تعاملی؛ رمزهای عبور هش‌شده؛ \lr{RBAC}؛ مقاومت در برابر \lr{SQL Injection} و \lr{XSS}} \\
\hline
\rl{کارایی} & \rl{زمان پاسخ‌گویی درخواست‌های معمولی کمتر از ۳ ثانیه؛ پشتیبانی از چندین کاربر همزمان بدون افت عملکرد} \\
\hline
\rl{دسترس‌پذیری} & \rl{در دسترس بودن ۲۴/۷؛ اختلال حداقلی در زمان به‌روزرسانی} \\
\hline
\rl{مقیاس‌پذیری} & \rl{امکان افزودن واحدها و کاربران جدید بدون تغییرات اساسی در سیستم} \\
\hline
\rl{قابلیت اطمینان} & \rl{ذخیره‌سازی پایدار اطلاعات؛ جلوگیری از از دست رفتن داده در صورت خطا؛ پشتیبان‌گیری دوره‌ای} \\
\hline
\rl{یکپارچگی داده} & \rl{جلوگیری از ثبت پرداخت تکراری؛ جلوگیری از رزروهای هم‌زمان و متداخل} \\
\hline
\rl{قابلیت استفاده} & \rl{رابط ساده برای همه‌ی نقش‌ها بدون نیاز به آموزش تخصصی؛ نمایش صحیح در مرورگرهای رایج} \\
\hline
\rl{قابلیت نگهداری} & \rl{کد ماژولار و قابل توسعه؛ مستندات فنی برای بخش‌های اصلی؛ امکان افزودن قابلیت‌های جدید بدون بازنویسی کامل} \\
\hline
\end{tabular}
\end{table}
